% !PRINTED 293
\weberpage{293}
%<<CONT
all the following groups $P_{\lambda+1}, P_{\lambda+2}, \ldots$ as well, which are indeed all
divisors of $P_\lambda$, are intransitive.

We must now recall the theorems 1., 2., 3. in \S.\,158 of the first volume. Since
$P_\lambda$ is a normal divisor of $P_{\lambda-1}$, it follows from those theorems
that $P_{\lambda-1}$ is imprimitive, and if by the necessary adjunctions the
group of $f(x) = 0$ has been reduced to $P_{\lambda-1}$, then by the further
adjunction of a root of a normal equation of degree $j_\lambda$ the function
$f(x)$ will split into several irreducible factors
\weq{2}{f(x) = f_1(x)\, f_2(x) \ldots f_h(x)}
which are all of equal degree and whose number $h$ is a divisor of
$j_\lambda$.

If we denote by $m$ the degree of $f(x)$, and by $m_1$ the degree of $f_1(x)$, then
\weq{3}{m = h\, m_1.}
By the cited theorem 1. in \S.\,158, Bd. I, the equations
$f_1 = 0, f_2 = 0, \ldots, f_h = 0$ all have the same group, which we shall denote
by $Q$. If $\alpha, \alpha_1, \alpha_2, \ldots, \alpha_{m_1-1}$ are the roots of
$f_1 = 0$, then we obtain the group $Q$ by collecting the permutations of the $\alpha$
that are brought about by $P_\lambda$. The first question concerns
the relation of the degrees $p_\lambda$ and $q$ of the two groups $P_\lambda$ and $Q$.

There may be several distinct permutations in $P_\lambda$ that generate the same
element in $Q$, that is, which contain the same permutation of the $\alpha$.
If $\pi_1, \pi_2$ are two distinct permutations from $P_\lambda$ that bring about
the same permutation among the $\alpha$, then $\pi_2 \pi_1^{-1} = \pi_0$ will be a
permutation that leaves the $\alpha$ at rest, and thus
$\pi_2 = \pi_0 \pi_1$. If conversely $\pi_0$ is any permutation from
$P_\lambda$ that does not permute the roots $\alpha$, then the permutation
$\pi_2 = \pi_0 \pi_1$ will produce the same change among the $\alpha$ as $\pi_1$.

From this it follows that every permutation of the $\alpha$ is generated equally often by the
permutations of $P_\lambda$, namely just as often as the $\alpha$
remain unchanged by permutations from $P_\lambda$. If we denote this number
by $p_0$, then
\weq{4}{p_\lambda = p_0\, q,}
\gesperrt{or the degree $p_\lambda$ of the group $P_\lambda$ is a multiple of the degree $q$
of $Q$.}
